\documentclass[12pt,a4paper]{article}
\usepackage[utf8]{inputenc}
\usepackage[vietnamese]{babel}
\usepackage{geometry}
\usepackage{amsmath}
\usepackage{amssymb}
\usepackage{graphicx}
\usepackage{booktabs}
\usepackage{longtable}
\usepackage{listings}
\usepackage{xcolor}
\usepackage{hyperref}
\usepackage{array}
\usepackage{float}

\geometry{left=2.5cm,right=2.5cm,top=2.5cm,bottom=2.5cm}

\definecolor{codegreen}{rgb}{0,0.6,0}
\definecolor{codegray}{rgb}{0.5,0.5,0.5}
\definecolor{codepurple}{rgb}{0.58,0,0.82}
\definecolor{backcolour}{rgb}{0.95,0.95,0.92}

\lstdefinestyle{mystyle}{
    backgroundcolor=\color{backcolour},
    commentstyle=\color{codegreen},
    keywordstyle=\color{blue},
    numberstyle=\tiny\color{codegray},
    stringstyle=\color{codepurple},
    basicstyle=\ttfamily\footnotesize,
    breaklines=true,
    keepspaces=true,
    numbers=left,
    numbersep=5pt,
    showspaces=false,
    showstringspaces=false,
    showtabs=false,
    tabsize=2,
    frame=single
}
\lstset{style=mystyle}

\title{\textbf{BÁO CÁO SESSION 10\\BREADTH-FIRST SEARCH TRONG SYNCHRONOUS NETWORK}}
\author{
  \textbf{Môn học:} Lập trình song song với CUDA\\
  \textbf{Nội dung:} BFS đồng bộ trên đồ thị vô hướng với CPU baseline và GPU CUDA
}
\date{\today}

\begin{document}

\maketitle
\tableofcontents
\newpage

\section{Giới thiệu}

\subsection{Bài toán}
Trong mô hình \textit{synchronous network}, thuật toán chạy theo các vòng (round) rời rạc.
Ở mỗi round, một tiến trình (đỉnh) có thể:
\begin{itemize}
    \item Gửi thông điệp thăm dò tới các đỉnh lân cận.
    \item Nhận thông điệp từ hàng xóm.
    \item Cập nhật trạng thái cục bộ (đã thăm/chưa thăm, mức BFS, cha trong cây BFS).
\end{itemize}

Bài toán BFS synchronous bắt đầu từ một đỉnh \texttt{start}, sau đó mở rộng theo từng lớp mức
(level) giống duyệt BFS cổ điển, nhưng nhấn mạnh khái niệm đồng bộ theo round.

\subsection{Mục tiêu thực hiện}
\begin{enumerate}
    \item Xây dựng chương trình BFS synchronous trên CPU (baseline).
    \item Xây dựng phiên bản GPU CUDA theo cơ chế frontier từng round.
    \item So sánh về:
    \begin{itemize}
        \item Thời gian thực thi.
        \item Số round giao tiếp.
        \item Tổng số thông điệp gửi trên cạnh có hướng.
        \item Độ đúng kết quả (mảng \texttt{level} CPU vs GPU).
    \end{itemize}
    \item Đóng gói quy trình chạy trong notebook Colab theo phong cách:
    \texttt{\%\%writefile} $\rightarrow$ \texttt{nvcc} $\rightarrow$ chạy chương trình.
\end{enumerate}

\section{Thiết kế dữ liệu và mô hình}

\subsection{Đầu vào JSON}
Đồ thị được nhập từ file JSON với định dạng:
\begin{lstlisting}[language={}, caption={Định dạng JSON đầu vào}]
{
  "nodes": [0, 1, 2, 3, 4],
  "edges": [[0,1], [0,2], [1,3], [2,4]]
}
\end{lstlisting}

Trong đó:
\begin{itemize}
    \item \texttt{nodes}: danh sách id đỉnh (có thể suy ra từ \texttt{edges}).
    \item \texttt{edges}: danh sách cạnh vô hướng \texttt{[u,v]}.
\end{itemize}

\subsection{Biểu diễn đồ thị trên GPU}
Đồ thị được chuyển sang định dạng CSR (Compressed Sparse Row):
\begin{itemize}
    \item \texttt{rowPtr}: mảng offset kích thước $n+1$.
    \item \texttt{colIdx}: danh sách đỉnh kề liên tục.
\end{itemize}

Ưu điểm của CSR:
\begin{itemize}
    \item Truy cập hàng xóm của mỗi đỉnh nhanh.
    \item Gọn bộ nhớ hơn ma trận kề với đồ thị thưa.
    \item Phù hợp với kernel dạng duyệt danh sách kề.
\end{itemize}

\section{Thuật toán BFS synchronous}

\subsection{CPU baseline}
Tại mỗi round:
\begin{enumerate}
    \item Duyệt toàn bộ đỉnh thuộc \texttt{frontier} hiện tại.
    \item Gửi thông điệp đến các hàng xóm.
    \item Nếu hàng xóm chưa thăm và chưa được lên lịch trong round này thì:
    \begin{itemize}
        \item Ghi \texttt{level = currentLevel + 1}
        \item Ghi \texttt{parent}
        \item Đưa vào \texttt{nextFrontier}
    \end{itemize}
\end{enumerate}

\subsection{GPU CUDA}
Mỗi thread xử lý một đỉnh \texttt{u}. Nếu \texttt{u} nằm trong frontier thì thread duyệt danh sách kề:
\begin{itemize}
    \item Tăng bộ đếm thông điệp bằng \texttt{atomicAdd(messageCount, 1)}.
    \item Dùng \texttt{atomicCAS} trên \texttt{visited[v]} để đảm bảo chỉ thread đầu tiên đánh dấu đỉnh \texttt{v}.
    \item Thread thắng quyền đánh dấu sẽ cập nhật \texttt{parent[v]}, \texttt{level[v]} và bật cờ \texttt{nextFrontier[v]}.
\end{itemize}

\subsection{Pseudo-code}
\begin{lstlisting}[language={}, caption={Pseudo-code BFS synchronous (frontier-based)}]
Initialize frontier = {start}
visited[start] = true
level[start] = 0
round = 0

while frontier not empty:
    nextFrontier = {}
    for each u in frontier (song song):
        for each v in Adj[u]:
            messages += 1
            if visited[v] == false:
                visited[v] = true
                parent[v] = u
                level[v] = round + 1
                add v to nextFrontier
    frontier = nextFrontier
    round += 1
\end{lstlisting}

\section{Notebook triển khai}
File notebook được tạo tại:
\begin{center}
\texttt{ss10/bfs\_sync\_cuda.ipynb}
\end{center}

Notebook gồm các phần chính:
\begin{enumerate}
    \item Kiểm tra GPU bằng \texttt{nvidia-smi}.
    \item Tạo \texttt{graph\_sample.json}.
    \item Tạo mã nguồn \texttt{bfs\_sync\_cuda.cu} bằng \texttt{\%\%writefile}.
    \item Biên dịch bằng:
    \begin{lstlisting}[language=bash]
!nvcc -O2 -o bfs_sync_cuda bfs_sync_cuda.cu -arch=sm_75
    \end{lstlisting}
    \item Chạy chương trình:
    \begin{lstlisting}[language=bash]
!./bfs_sync_cuda graph_sample.json 0
    \end{lstlisting}
    \item Chạy benchmark nhanh với nhiều \texttt{start node}.
\end{enumerate}

\section{Kết quả và phân tích}

\subsection{Các chỉ số theo dõi}
Chương trình in ra:
\begin{itemize}
    \item \texttt{time(ms)} của CPU và GPU.
    \item \texttt{communication rounds}.
    \item \texttt{messages} (số lần duyệt cạnh có hướng).
    \item Bảng \texttt{node\_id, cpu\_level, gpu\_level, gpu\_parent}.
    \item Round trace của GPU.
\end{itemize}

\subsection{Nhận xét tổng quan}
\begin{itemize}
    \item Với đồ thị nhỏ, thời gian GPU có thể chưa vượt trội do overhead launch kernel và sao chép dữ liệu.
    \item Với đồ thị lớn và frontier rộng, GPU có lợi thế nhờ xử lý song song số lượng lớn cạnh trong mỗi round.
    \item Việc dùng \texttt{atomicCAS} giúp đảm bảo đúng tính chất BFS theo lớp khi nhiều thread tranh chấp cùng một đỉnh chưa thăm.
\end{itemize}

\section{Hạn chế và hướng cải tiến}

\subsection{Hạn chế hiện tại}
\begin{itemize}
    \item Parser JSON trong mã CUDA dùng regex đơn giản, phù hợp cho dữ liệu có cấu trúc chuẩn.
    \item Mỗi round có đồng bộ host-device để lấy \texttt{nextFrontier} và đếm số đỉnh mới.
    \item Chưa tối ưu cho đồ thị rất lớn bằng các kỹ thuật nén frontier nâng cao.
\end{itemize}

\subsection{Hướng mở rộng}
\begin{enumerate}
    \item Dùng định dạng input binary/CSR trực tiếp để giảm chi phí parse.
    \item Tối ưu frontier bằng bitset và kỹ thuật \textit{direction-optimizing BFS}.
    \item Dùng nhiều benchmark dataset (road network, social network) để đánh giá toàn diện.
    \item Kết nối kết quả CUDA với ứng dụng .NET UI để so sánh 4 chế độ: Thread/Task, Parallel.For, Channel, CUDA.
\end{enumerate}

\section{Kết luận}
Báo cáo đã trình bày cách hiện thực BFS trong synchronous network bằng CUDA, với quy trình chạy rõ ràng trên notebook.
Cấu trúc triển khai bám sát mô hình round-based của bài giảng, đồng thời có đối chiếu CPU để xác thực.
Kết quả cho thấy phương pháp GPU có tiềm năng tăng tốc khi kích thước đồ thị đủ lớn và mức song song cao.

\section{Tài liệu tham khảo}
\begin{enumerate}
    \item NVIDIA Corporation, \textit{CUDA C++ Programming Guide}.
    \item Mark Harris, \textit{Optimizing Parallel Reduction in CUDA}.
    \item Tài liệu bài giảng: \textit{Breadth-First Search in Synchronous Networks}.
\end{enumerate}

\end{document}
